\section{Grassmann--Cayley extensor algebra}
\label{sec:molecular}

This chapter opens the \textbf{molecular-conjecture program}
(Phases~17--26), the project's longest-horizon target: the molecular
conjecture of Tay and Whiteley, proved by
Katoh--Tanigawa~\cite{katohTanigawa2011}. The program-level plan ---
target, five-strata architecture, and the phase breakdown --- is
described in the project's \texttt{notes/MolecularConjecture.md}. The
program is authored one chapter per phase: this chapter is the
forward-mode authoritative dep-graph and lemma index for Phase~17 (the
extensor algebra,
\cref{sec:molecular-homog,sec:molecular-extensor,sec:molecular-lemma21}),
and the panel-hinge rigidity matrix of Phase~18 follows in
\cref{sec:molecular-rigidity-matrix}.

This chapter formalizes the \emph{Grassmann--Cayley / extensor algebra}
layer of Katoh--Tanigawa~\cite[\S2.1]{katohTanigawa2011}, the first of
the five strata. It is genuinely new linear algebra --- mathlib carries
the exterior algebra $\bigwedge[\R] M$ but not the Plücker/extensor
coordinatization, the join, or the affine-subspace--to-extensor map ---
and bottoms out the most difficult case of the conjecture's algebraic
induction (Case~III, Phases~22--23). The load-bearing target is
\textbf{Lemma~2.1} (\cref{lem:extensor-independence}): the
$D = \binom{d+1}{2}$ many $(d-1)$-extensors associated to $d+1$ affinely
independent points are linearly independent.

\medskip
\noindent\emph{Forward-mode chapter.} Following the convention for
Phases~6 onward, this chapter is authored as the phase's authoritative
dep-graph: each node carries its intended \texttt{\textbackslash uses\{\}}
dependency structure, and turns green as its Lean lands. The Phase~17
nodes (\cref{sec:molecular-homog,sec:molecular-extensor,sec:molecular-lemma21})
are all formalized
(\texttt{\textbackslash lean\{\}} / \texttt{\textbackslash leanok}); the
dep-graph there records the realized proof structure of
\texttt{Molecular/Extensor.lean}.

\subsection{Homogeneous coordinates and affine independence}
\label{sec:molecular-homog}

We work in $\R^{d}$ with points homogenized into $\R^{d+1}$. Throughout,
$D = \binom{d+1}{2} = \tfrac{d(d+1)}{2}$ is the dimension of the space of
infinitesimal screw motions of a rigid body in $\R^{d}$, matching the
$\delta = \mathrm{bodyBarDim}\,n$ of the body-bar / body-hinge chapters
(\cref{sec:body-bar,sec:body-hinge}).

\begin{definition}[Homogeneous coordinatization]
  \label{def:homogeneous-coords}
  \lean{CombinatorialRigidity.Molecular.homogenize}
  \leanok
  For a point $p \in \R^{d}$ write $\bar p = (p, 1) \in \R^{d+1}$ for its
  homogeneous coordinate vector. A finite family $p_1, \dots, p_k$ of
  points of $\R^{d}$ is \emph{affinely independent} when the homogenized
  family $\bar p_1, \dots, \bar p_k$ is linearly independent in
  $\R^{d+1}$.
\end{definition}

\begin{lemma}[Affine independence via homogenization and the top extensor]
  \label{lem:affine-indep-iff}
  \lean{CombinatorialRigidity.Molecular.affineIndependent_iff_linearIndependent_homogenize,
    CombinatorialRigidity.Molecular.affineIndependent_fin_iff_det_homogenize}
  \leanok
  \uses{def:homogeneous-coords}
  Points $p_1, \dots, p_k \in \R^{d}$ are affinely independent if and
  only if the homogenized family $\bar p_1, \dots, \bar p_k$ is linearly
  independent in $\R^{d+1}$ --- equivalently, once the join is in hand
  (\cref{def:join}), if and only if the join
  $\bar p_1 \vee \cdots \vee \bar p_k$ is nonzero. In particular $d+1$
  points are affinely independent iff their top extensor
  $\bar p_1 \vee \cdots \vee \bar p_{d+1}$ --- the determinant of the
  $(d+1)\times(d+1)$ matrix of their homogeneous coordinates --- is
  nonzero.
\end{lemma}
\begin{proof}
  \leanok
  Affine independence of $p_1, \dots, p_k$ means every coefficient
  family $w$ with $\sum_i w_i = 0$ and $\sum_i w_i p_i = 0$ vanishes; a
  vanishing combination $\sum_i w_i \bar p_i = 0$ of the homogenized
  vectors is exactly such a family (the last coordinate gives
  $\sum_i w_i = 0$, the first $d$ give $\sum_i w_i p_i = 0$), so affine
  independence is linear independence of the $\bar p_i$. Once the join is
  available it is the image in $\bigwedge^{k} \R^{d+1}$, nonzero exactly
  when the vectors are linearly independent. For $k = d+1$ the top
  exterior power is one-dimensional and the coordinate of the join is the
  determinant of the homogeneous coordinate matrix; that determinant is
  nonzero iff the rows are linearly independent.
\end{proof}

\subsection{Extensors, the join, and Plücker coordinates}
\label{sec:molecular-extensor}

\begin{definition}[Extensor]
  \label{def:extensor}
  \lean{CombinatorialRigidity.Molecular.extensor}
  \leanok
  A \emph{$j$-extensor} in $\R^{d+1}$ is a decomposable element of the
  $j$-th exterior power $\bigwedge^{j} \R^{d+1}$, i.e.\ one of the form
  $v_1 \wedge \cdots \wedge v_j$ for vectors $v_i \in \R^{d+1}$
  (Barnabei--Brini--Rota~\cite{barnabeiBriniRota1985}). We build the
  symbolic layer on mathlib's $\bigwedge[\R] \R^{d+1}$ --- the extensor is
  the value of mathlib's alternating multilinear map into the exterior
  algebra, landing in the graded piece $\bigwedge^{j} \R^{d+1}$ --- and
  recover the geometric notion of the
  extensor of an affine subspace below. Being alternating, the extensor
  vanishes whenever two of its vectors coincide (the load-bearing fact for
  \cref{lem:extensor-independence}).
\end{definition}

\begin{definition}[Join]
  \label{def:join}
  \lean{CombinatorialRigidity.Molecular.join}
  \leanok
  \uses{def:extensor}
  The \emph{join} $A \vee B$ of extensors is their exterior product
  $A \wedge B$ (the span of the corresponding subspaces, when they meet
  only at the origin). On extensors it satisfies
  $(\bar v_1 \vee \cdots \vee \bar v_m) \vee (\bar w_1 \vee \cdots \vee \bar w_n)
   = \bar v_1 \vee \cdots \vee \bar v_m \vee \bar w_1 \vee \cdots \vee \bar w_n$,
  the extensor of the concatenated family. It is alternating and vanishes
  whenever a vector is repeated (inherited from \cref{def:extensor}), and is
  associative and graded-commutative as inherited from the exterior product.
\end{definition}

\begin{definition}[Plücker coordinates]
  \label{def:plucker-coords}
  \lean{CombinatorialRigidity.Molecular.pluckerCoord,
    CombinatorialRigidity.Molecular.pluckerVector}
  \leanok
  \uses{def:extensor, def:join}
  Assemble a family $v_1, \dots, v_j \in \R^{d+1}$ into the
  $j \times (d+1)$ matrix $A(v)$ whose $i$-th row is $v_i$. The
  \emph{Plücker coordinate} $P_{i_1, \dots, i_j}$ of the $j$-extensor
  $v_1 \wedge \cdots \wedge v_j$, indexed by an increasing column set
  $1 \le i_1 < \cdots < i_j \le d+1$, is the
  Katoh--Tanigawa~\cite[\S2.1]{katohTanigawa2011} sign
  $(-1)^{1 + \sum_\ell i_\ell}$ times the determinant of the
  $j \times j$ submatrix of $A(v)$ on those columns. The \emph{Plücker
  coordinate vector} collects these over all $\binom{d+1}{j}$ column
  sets. This is the coordinatized bridge from the symbolic
  exterior-algebra layer to the concrete minor vectors on which
  Lemma~2.1 and the downstream rank computations operate; for the top
  case $j = d+1$ the single coordinate is, up to the sign, the full
  homogeneous-coordinate determinant of \cref{lem:affine-indep-iff}.
\end{definition}

\begin{definition}[Extensor of an affine subspace]
  \label{def:affine-subspace-extensor}
  \lean{CombinatorialRigidity.Molecular.affineSubspaceExtensor,
    CombinatorialRigidity.Molecular.affineSubspaceExtensor_ne_zero_iff}
  \leanok
  \uses{def:homogeneous-coords, def:extensor, def:join}
  For an affine subspace spanned by affinely independent points
  $p_1, \dots, p_k \in \R^{d}$, its \emph{extensor} is
  $C(\cdot) = \bar p_1 \vee \cdots \vee \bar p_k \in \bigwedge^{k} \R^{d+1}$,
  well-defined up to a nonzero scalar. By \cref{lem:affine-indep-iff} it
  is nonzero exactly when the points are affinely independent, so $C$
  faithfully encodes the affine subspace. (Formally, the nonvanishing
  characterization combines the homogenization bridge of
  \cref{lem:affine-indep-iff} with the fact that a $k$-extensor is nonzero
  iff its vectors are linearly independent --- the converse direction
  extends the family to a basis of its span and reads off a nonzero
  exterior-power basis coordinate.)
\end{definition}

\subsection{Independence of the $(d-1)$-extensors (Lemma 2.1)}
\label{sec:molecular-lemma21}

\begin{lemma}[Extensor independence; Katoh--Tanigawa Lemma 2.1]
  \label{lem:extensor-independence}
  \lean{CombinatorialRigidity.Molecular.omitTwoExtensor,
    CombinatorialRigidity.Molecular.omitTwoExtensor_linearIndependent,
    CombinatorialRigidity.Molecular.omitTwoExtensor_linearIndependent_of_li}
  \leanok
  \uses{def:affine-subspace-extensor, def:plucker-coords, lem:affine-indep-iff}
  Let $p_1, \dots, p_{d+1} \in \R^{d}$ be affinely independent. The
  $D = \binom{d+1}{2}$ many $(d-1)$-extensors obtained by omitting two of
  the $d+1$ points,
  \[
    \bar p_1 \vee \cdots \vee \widehat{\bar p_a} \vee \cdots
      \vee \widehat{\bar p_b} \vee \cdots \vee \bar p_{d+1}
    \qquad (1 \le a < b \le d+1),
  \]
  are linearly independent in $\bigwedge^{d-1} \R^{d+1}$.
\end{lemma}
\begin{proof}
  \leanok
  \uses{def:join, lem:affine-indep-iff}
  Suppose a linear combination of the $D$ extensors vanishes. Fix a pair
  $\{a, b\}$ and join the relation on the left with the $2$-extensor
  $\bar p_a \vee \bar p_b$. Every term whose omitted pair is not
  $\{a, b\}$ then repeats a vector $\bar p_a$ or $\bar p_b$ (one of $a, b$
  survives the omission of the other pair), so by the alternating property
  of the join (\cref{def:join}) it vanishes; the single surviving term is a
  nonzero scalar times the top extensor $\bar p_1 \vee \cdots \vee \bar
  p_{d+1}$, which is nonzero by affine independence
  (\cref{lem:affine-indep-iff}). Hence the coefficient on the
  $\{a, b\}$-extensor is zero. As $\{a, b\}$ was arbitrary, all
  coefficients vanish. The formalization reparametrizes $d = e + 1$ (so
  the $e + 2$ points are in $\R^{e+1}$ and each omit-two extensor is an
  $e$-extensor, clearing the $\mathbb{N}$-subtraction $d - 1$), indexes
  the $D$ extensors by the increasing pairs $a < b$, and reindexes the
  surviving join via the bijection that sends the first two slots to
  $a, b$ and the rest to the increasing enumeration of $\{a, b\}^{c}$.
  The argument reads only \emph{linear} independence of the
  homogenizations $\bar p_i$ (affine independence enters once, via
  \cref{lem:affine-indep-iff}, to supply it), so it holds verbatim for
  any linearly-independent family of $d + 1$ vectors --- the form the
  Case~III spanning step (\cref{lem:case-III-claim612-extensor-span})
  consumes.
\end{proof}

This independence is the linear-algebraic foundation that Case~III of
the conjecture's algebraic induction (Phases~22--23) bottoms out on: a
nonzero screw $r \in \R^{D}$ orthogonal to the extensors of $d+1$
distinct generic panels would, by \cref{lem:extensor-independence}, be
orthogonal to a spanning set of $\R^{D}$ and hence zero. The genuine
panel-hinge rigidity matrix that consumes this fact is built next, in
\cref{sec:molecular-rigidity-matrix}.
